\documentclass[border=4pt]{standalone}
\usepackage{amsmath,amssymb,mathtools,xcolor,varwidth}
\definecolor{motionblue}{HTML}{1E6FE0}
\definecolor{dirgreen}{HTML}{00A35A}
\definecolor{forceorange}{HTML}{E8770A}
\begin{document}
\begin{varwidth}{28em}
\large\bfseries
The \textcolor{motionblue}{change of motion} is made \textcolor{dirgreen}{in the direction of the right line} in which \textcolor{forceorange}{the force} is impressed. This motion, being always directed the same way with \textcolor{forceorange}{the generating force}, if the body moved before, is \textcolor{dirgreen}{added to or subducted from} the former motion — according as they directly conspire with or are directly contrary to each other; or \textcolor{dirgreen}{obliquely joined}, when they are oblique, so as to produce a new motion compounded from the determination of both. Thus \textcolor{motionblue}{$\Delta(\text{motion})$} runs \textcolor{dirgreen}{parallel} to \textcolor{forceorange}{$\mathbf{F}$}.
\end{varwidth}
\end{document}
